\chapter{Web の世界地図}
\sheet{1-1〜1-3}

道具をそろえる前に、これから作るものの全体像を見ておきます。
細かい言葉は覚えなくて大丈夫です。この章は、迷ったときに戻ってくる「地図」として使ってください。

\section{Web ページが表示されるまで}

ブラウザに \code{https://rikotaikendo.jimdofree.com} と打ち込んでから、ページが表示されるまでに起きていることは、次の図の通りです。

\begin{center}
\begin{tikzpicture}[font=\sffamily\small,
  box/.style={draw=kon,rounded corners=3pt,minimum width=3.2cm,minimum height=2.2cm,align=center,fill=white},
  arr/.style={-{Stealth[length=3mm]},thick}]
  \node[box] (b) {\faLaptop\\[2pt]\textbf{ブラウザ}\\{\scriptsize あなたの PC}};
  \node[box,right=6.6cm of b] (s) {\faServer\\[2pt]\textbf{サーバー}\\{\scriptsize インターネットの向こう}};
  \draw[arr,kon] ([yshift=4mm]b.east) -- node[above,font=\scriptsize]{① ページをください（リクエスト）} ([yshift=4mm]s.west);
  \draw[arr,shu] ([yshift=-4mm]s.west) -- node[below,font=\scriptsize,align=center]{② ファイル一式を返す（レスポンス）} ([yshift=-4mm]b.east);
  \node[below=4mm of b,font=\scriptsize,align=center,text width=4cm] {③ 受け取ったファイルを組み立てて画面に描く};
\end{tikzpicture}
\end{center}

\begin{itemize}
  \item \term{ブラウザ}：Chrome や Edge のこと。受け取ったファイルを読んで画面を描きます。
  \item \term{サーバー}：ページの元になるファイルを置いておき、頼まれたら返すコンピュータです。
  \item \term{リクエスト}と\term{レスポンス}：「ください」と「どうぞ」のやりとりです。Web の通信はすべてこれの繰り返しです。
\end{itemize}

ブラウザが受け取るファイルは、主に次の3種類です。

\begin{center}\small
\begin{tabularx}{\linewidth}{l l X}
\toprule
ファイル & 役割 & たとえると \\
\midrule
\textbf{HTML} & 何が書いてあるか（文章・見出し・画像の場所） & 家の骨組み \\
\textbf{CSS} & どう見せるか（色・大きさ・配置） & 内装・ペンキ \\
\textbf{JavaScript} & どう動くか（クリックしたら開く、絞り込む） & 電気・スイッチ \\
\bottomrule
\end{tabularx}
\end{center}

第1巻の第7章で HTML と CSS を、第8章で JavaScript を書きます。

\section{フロントエンド・バックエンド・データベース}

部のサイトには、ニュースや大会結果を部員が更新できる「管理画面」があります。
更新された内容はどこかに保存しておかなければなりません。
保存先が\term{データベース}（DB）です。

\begin{center}
\begin{tikzpicture}[font=\sffamily\small,
  layer/.style={draw,rounded corners=3pt,minimum width=3.6cm,minimum height=1.6cm,align=center},
  arr/.style={{Stealth[length=2.5mm]}-{Stealth[length=2.5mm]},thick,gray}]
  \node[layer,draw=kon,fill=kon!6] (f) {\textbf{フロントエンド}\\{\scriptsize ブラウザで動く部分}\\{\scriptsize 見た目・操作}};
  \node[layer,draw=kincha,fill=kincha!8,right=1.6cm of f] (b) {\textbf{バックエンド}\\{\scriptsize サーバーで動く部分}\\{\scriptsize 計算・権限チェック}};
  \node[layer,draw=matcha,fill=matcha!8,right=1.6cm of b] (d) {\textbf{データベース}\\{\scriptsize データの保管庫}\\{\scriptsize 部員・ニュース・結果}};
  \draw[arr] (f)--(b); \draw[arr] (b)--(d);
  \node[below=2mm of f,font=\scriptsize] {第1〜2巻};
  \node[below=2mm of b,font=\scriptsize] {第2〜4巻};
  \node[below=2mm of d,font=\scriptsize] {第3〜4巻};
\end{tikzpicture}
\end{center}

\begin{tip}[なぜバックエンドを挟むのか]
ブラウザから直接データベースを書き換えられると、誰でも部員名簿を消せてしまいます。
そこで「この人は管理者か？」を確かめる係が必要になります。それがバックエンドの大事な仕事です。
第4巻で、この「確かめる係」を自分で作ります。
\end{tip}

\section{部のサイトの構成}

本番の \code{rikotaikendo-web} は、次のような部品の組み合わせでできています。
すべてを今理解する必要はありません。第6巻を終えたとき、この図のどの部品についても説明できるようになっていればゴールです。

\begin{center}
\begin{tikzpicture}[font=\sffamily\scriptsize,
  svc/.style={draw,rounded corners=3pt,minimum width=2.7cm,minimum height=1.15cm,align=center,fill=white},
  arr/.style={-{Stealth[length=2mm]},thick,gray!70}]
  \node[svc,draw=gray] (user) {\faUsers\ 訪問者\\新入生・部員・OB};
  \node[svc,draw=gray,below=1.2cm of user] (admin) {\faUserCog\ 管理者\\主務など};
  \node[svc,draw=kon,very thick,right=1.4cm of $(user)!0.5!(admin)$,minimum height=2.6cm,minimum width=3.2cm] (next)
    {\textbf{Next.js アプリ}\\[2pt]公開ページ\\管理画面\\[2pt]{\color{gray}置き場所：Vercel}};
  \node[svc,draw=matcha,right=1.4cm of next,yshift=1.3cm] (db) {\textbf{Supabase}\\DB・ログイン・画像};
  \node[svc,draw=kincha,right=1.4cm of next] (mail) {\textbf{Resend}\\問い合わせメール};
  \node[svc,draw=shu,right=1.4cm of next,yshift=-1.3cm] (google) {\textbf{Google}\\ログイン（OAuth）};
  \node[svc,draw=gray,above=0.9cm of next] (gh) {\faGithub\ \textbf{GitHub}\\コードの保管庫};
  \draw[arr] (user) -- (next); \draw[arr] (admin) -- (next);
  \draw[arr] (next) -- (db); \draw[arr] (next) -- (mail); \draw[arr] (next) -- (google);
  \draw[arr,dashed] (gh) -- node[right]{push すると自動で公開} (next);
\end{tikzpicture}
\end{center}

\begin{center}\small
\begin{tabularx}{\linewidth}{l X l}
\toprule
部品 & 何をしているか & 学ぶ巻 \\
\midrule
GitHub & コードを保管し、変更の履歴を残す。チームでの共同作業の場 & 第1巻〜 \\
Next.js & ページを作る仕組み（フレームワーク） & 第2巻 \\
Vercel & Next.js のアプリをインターネットに公開する場所 & 第3巻 \\
Supabase & データベース、ログイン、画像の保存をまとめて提供するサービス & 第3〜4巻 \\
Google OAuth & 「Google でログイン」の仕組み & 第4巻 \\
Resend & 問い合わせフォームからメールを送るサービス & 第5巻 \\
\bottomrule
\end{tabularx}
\end{center}

\section{字幕アプリの構成}

すでに運用で使っている字幕アプリも、ほとんど同じ部品でできています。
違うのは、\textbf{重い動画処理を GitHub Actions に任せている}点です。

\begin{center}
\begin{tikzpicture}[font=\sffamily\scriptsize,
  st/.style={draw=kon,rounded corners=3pt,minimum width=2.15cm,minimum height=1.25cm,align=center,fill=white,text width=2.05cm},
  arr/.style={-{Stealth[length=2mm]},thick,kon!70}]
  \node[st] (a) {① 試合情報を入力して動画を選ぶ\\{\color{gray}ブラウザ}};
  \node[st,right=4mm of a] (b) {② 動画を一時保存\\{\color{gray}Google Drive}};
  \node[st,right=4mm of b] (c) {③ 処理を依頼\\{\color{gray}GitHub Actions}};
  \node[st,right=4mm of c] (d) {④ 字幕を焼き込む\\{\color{gray}FFmpeg}};
  \node[st,below=8mm of d] (e) {⑤ 完了を知らせる\\{\color{gray}Webhook}};
  \node[st,left=4mm of e] (f) {⑥ 状態を「完了」に\\{\color{gray}Supabase}};
  \node[st,left=4mm of f] (g) {⑦ YouTube に限定公開\\{\color{gray}Actions}};
  \node[st,left=4mm of g] (h) {⑧ 一定時間後に動画を削除\\{\color{gray}定期実行}};
  \draw[arr] (a)--(b); \draw[arr] (b)--(c); \draw[arr] (c)--(d); \draw[arr] (d)--(e);
  \draw[arr] (e)--(f); \draw[arr] (f)--(g); \draw[arr] (g)--(h);
\end{tikzpicture}
\end{center}

10月の大会では、外部サーバーの不調で3本の処理が失敗しました。
あれは、この図の③〜④のどこかで止まっていたということです。
第6巻では、こうした失敗が「どこで」起きたのかを自分で突き止められるようにします。

\begin{step}{三角と一緒に、字幕アプリの裏側をたどる（定例で実施）}
\sheet{1-3}
定例の時間を1回使い、三角の画面共有で、試合動画1本の処理を最初から最後まで追いかけます。
見ながら、下の表の空欄を埋めてください。

\medskip
\begin{tabularx}{\linewidth}{l X}
\toprule
上の図の番号 & どの画面で確認できたか／何が表示されていたか \\
\midrule
① 入力 & \\[6pt]
③ Actions の実行 & \\[6pt]
④ FFmpeg のログ & \\[6pt]
⑥ DB の状態 & \\[6pt]
⑦ YouTube & \\[6pt]
\bottomrule
\end{tabularx}
\end{step}

\section{ローカルとクラウド}

最後に、言葉を2つだけ覚えておきます。

\begin{itemize}
  \item \term{ローカル}：自分の PC の中のこと。作業中のファイルはまずここにあります。
  \item \term{クラウド}（リモート）：インターネットの向こうにあるコンピュータのこと。GitHub や Vercel がこれにあたります。
\end{itemize}

開発は「ローカルで作って試す → クラウドに送って公開する」の繰り返しです。
ローカルから GitHub に送る操作を \term{push}（プッシュ）と呼びます。第5章で実際にやります。

\begin{checklist}
\begin{itemize}
  \checkitem ブラウザとサーバーのやりとりを、リクエスト・レスポンスという言葉を使って説明できる
  \checkitem HTML・CSS・JavaScript の役割の違いを言える
  \checkitem 管理画面とデータベースの間に、なぜバックエンドが必要なのかを言える
  \checkitem 字幕アプリの処理が、どの順番でどのサービスを通るのかを図で示せる
\end{itemize}
\end{checklist}
